\chapter{Experimental Evaluation}

This chapter reports the experiments in the order of the research questions. It first evaluates lexical and hybrid retrieval, then isolates the incremental contribution of E5, then examines BCA coverage and component-aware reliability, and finally analyzes legal-domain adaptation. The central adaptation finding is that dense retriever improvement does not necessarily transfer into a complete hybrid retrieval system.

\section{Datasets and Experimental Setup}

Quantitative table values are displayed to four decimal places, using scientific notation with four fractional digits in the mantissa for tiny nonzero values. Deltas are calculated before display rounding; integer counts and indices remain integers.

The evaluated conditions and corpus sizes are defined in Table~\ref{tab:4.1}. ALQAC and Zalo each have clean and parsed corpus representations; their A/B conditions share query identities and are therefore paired corpus conditions rather than independent datasets. BCA is evaluated separately because many questions share legal targets and because some targets are unavailable in the indexed corpus.

NDCG@10 is the primary endpoint. Hit@10, MRR@10, and MAP@10 are reported to distinguish top-ten recovery, first-relevant ordering, and multiple-target ranking. ALQAC and Zalo use query-level paired comparisons for the two-way versus three-way hybrid. BCA uses component-aware analysis in addition to descriptive query-level summaries. These units of inference are not interchangeable.

\begingroup
\small
\setlength{\tabcolsep}{3pt}
\renewcommand{\arraystretch}{1.18}
\begin{longtable}{>{\raggedright\arraybackslash}p{2.938cm}>{\raggedright\arraybackslash}p{6.569cm}>{\raggedright\arraybackslash}p{5.833cm}}
\caption[System identities and comparison boundaries.]{System identities and comparison boundaries.}\label{tab:5.1}\\
\toprule
\textbf{Configuration} & \textbf{Baseline role} & \textbf{Adaptation role} \\
\midrule
\endfirsthead
\toprule
\textbf{Configuration} & \textbf{Baseline role} & \textbf{Adaptation role} \\
\midrule
\endhead
\midrule
\multicolumn{3}{r}{\footnotesize Continued on next page}\\
\endfoot
\bottomrule
\endlastfoot
BM25 & Fixed lexical reference & Unchanged lexical component \\
BGE-M3 dense & Pretrained dense retriever & Adapted in the legal-domain fine-tuning studies \\
E5 dense & Additional pretrained dense retriever & Frozen throughout the thesis \\
BM25+\allowbreak{}BGE-M3 & Two-way lexical–semantic hybrid & Evaluated with pretrained and adapted BGE-M3 \\
BM25+\allowbreak{}BGE-M3+\allowbreak{}E5 & Three-way hybrid & Evaluated with pretrained and adapted BGE-M3 while E5 remains fixed \\
\end{longtable}
\endgroup

The experiments compare methods within their own query populations and selection boundaries. Results from different folds, corpora, or adaptation recipes are not pooled into a single cross-protocol score.

\section{Lexical and Hybrid Retrieval}

This section compares lexical–semantic hybrid retrieval with the fixed BM25 reference. Table~\ref{tab:5.2} reports BM25, pretrained dense retrievers, and hybrid systems for ALQAC and Zalo. Table~\ref{tab:5.3} summarizes the hybrid changes relative to BM25 and includes BCA's combined descriptive view.

\begingroup
\footnotesize
\setlength{\tabcolsep}{3pt}
\renewcommand{\arraystretch}{1.18}
\begin{longtable}{>{\raggedright\arraybackslash}p{2.363cm}>{\raggedright\arraybackslash}p{2.865cm}>{\raggedright\arraybackslash}p{2.363cm}>{\raggedright\arraybackslash}p{2.363cm}>{\raggedright\arraybackslash}p{2.363cm}>{\raggedright\arraybackslash}p{2.363cm}}
\caption[Pretrained standalone and hybrid results, query-macro OOF population.]{Pretrained standalone and hybrid results, query-macro OOF population.}\label{tab:5.2}\\
\toprule
\textbf{Condition} & \textbf{Method} & \textbf{NDCG@10} & \textbf{MRR@10} & \textbf{MAP@10} & \textbf{Hit@10} \\
\midrule
\endfirsthead
\toprule
\textbf{Condition} & \textbf{Method} & \textbf{NDCG@10} & \textbf{MRR@10} & \textbf{MAP@10} & \textbf{Hit@10} \\
\midrule
\endhead
\midrule
\multicolumn{6}{r}{\footnotesize Continued on next page}\\
\endfoot
\bottomrule
\endlastfoot
ALQAC A & BM25 & 0.7986 & 0.7648 & 0.7621 & 0.9112 \\
ALQAC A & BGE-M3 pretrained & 0.8437 & 0.8134 & 0.8108 & 0.9433 \\
ALQAC A & E5 pretrained & 0.8126 & 0.7729 & 0.7722 & 0.9371 \\
ALQAC A & BM25+\allowbreak{}BGE-M3 & 0.8805 & 0.8534 & 0.8516 & 0.9692 \\
ALQAC A & BM25+\allowbreak{}BGE-M3+\allowbreak{}E5 & 0.8784 & 0.8519 & 0.8504 & 0.9630 \\
ALQAC B & BM25 & 0.6479 & 0.6175 & 0.6159 & 0.7472 \\
ALQAC B & BGE-M3 pretrained & 0.6684 & 0.6401 & 0.6382 & 0.7620 \\
ALQAC B & E5 pretrained & 0.6510 & 0.6177 & 0.6165 & 0.7583 \\
ALQAC B & BM25+\allowbreak{}BGE-M3 & 0.7067 & 0.6812 & 0.6797 & 0.7904 \\
ALQAC B & BM25+\allowbreak{}BGE-M3+\allowbreak{}E5 & 0.7055 & 0.6805 & 0.6793 & 0.7867 \\
Zalo A & BM25 & 0.5264 & 0.4731 & 0.4682 & 0.7100 \\
Zalo A & BGE-M3 pretrained & 0.6424 & 0.5917 & 0.5850 & 0.8235 \\
Zalo A & E5 pretrained & 0.5640 & 0.5104 & 0.5046 & 0.7509 \\
Zalo A & BM25+\allowbreak{}BGE-M3 & 0.6805 & 0.6298 & 0.6235 & 0.8583 \\
Zalo A & BM25+\allowbreak{}BGE-M3+\allowbreak{}E5 & 0.6818 & 0.6317 & 0.6254 & 0.8576 \\
Zalo B & BM25 & 0.5620 & 0.5099 & 0.5050 & 0.7406 \\
Zalo B & BGE-M3 pretrained & 0.6889 & 0.6431 & 0.6366 & 0.8526 \\
Zalo B & E5 pretrained & 0.6259 & 0.5748 & 0.5689 & 0.8054 \\
Zalo B & BM25+\allowbreak{}BGE-M3 & 0.7214 & 0.6767 & 0.6705 & 0.8789 \\
Zalo B & BM25+\allowbreak{}BGE-M3+\allowbreak{}E5 & 0.7344 & 0.6913 & 0.6845 & 0.8889 \\
\end{longtable}
\endgroup

BGE-M3 is the strongest individual retriever by NDCG@10 in all four ALQAC/Zalo conditions. Both hybrids exceed BM25, and both also exceed the strongest individual component in these conditions. The result supports lexical–semantic complementarity at the system level, although it does not identify the linguistic mechanism behind each gain.

\begingroup
\footnotesize
\setlength{\tabcolsep}{3pt}
\renewcommand{\arraystretch}{1.18}
\begin{longtable}{>{\raggedright\arraybackslash}p{3.075cm}>{\raggedright\arraybackslash}p{2.349cm}>{\raggedright\arraybackslash}p{2.197cm}>{\raggedright\arraybackslash}p{2.197cm}>{\raggedright\arraybackslash}p{2.431cm}>{\raggedright\arraybackslash}p{2.431cm}}
\caption[Hybrid NDCG@10 changes relative to fixed BM25.]{Hybrid NDCG@10 changes relative to fixed BM25.}\label{tab:5.3}\\
\toprule
\textbf{Condition} & \textbf{Hybrid} & \textbf{BM25} & \textbf{Hybrid score} & \textbf{Absolute change} & \textbf{Relative change} \\
\midrule
\endfirsthead
\toprule
\textbf{Condition} & \textbf{Hybrid} & \textbf{BM25} & \textbf{Hybrid score} & \textbf{Absolute change} & \textbf{Relative change} \\
\midrule
\endhead
\midrule
\multicolumn{6}{r}{\footnotesize Continued on next page}\\
\endfoot
\bottomrule
\endlastfoot
ALQAC A & BM25+\allowbreak{}BGE-M3 & 0.7986 & 0.8805 & +\allowbreak{}0.0819 & +\allowbreak{}10.2590\% \\
ALQAC A & BM25+\allowbreak{}BGE-M3+\allowbreak{}E5 & 0.7986 & 0.8784 & +\allowbreak{}0.0798 & +\allowbreak{}9.9970\% \\
ALQAC B & BM25+\allowbreak{}BGE-M3 & 0.6479 & 0.7067 & +\allowbreak{}0.0588 & +\allowbreak{}9.0690\% \\
ALQAC B & BM25+\allowbreak{}BGE-M3+\allowbreak{}E5 & 0.6479 & 0.7055 & +\allowbreak{}0.0576 & +\allowbreak{}8.8910\% \\
Zalo A & BM25+\allowbreak{}BGE-M3 & 0.5264 & 0.6805 & +\allowbreak{}0.1541 & +\allowbreak{}29.2760\% \\
Zalo A & BM25+\allowbreak{}BGE-M3+\allowbreak{}E5 & 0.5264 & 0.6818 & +\allowbreak{}0.1554 & +\allowbreak{}29.5210\% \\
Zalo B & BM25+\allowbreak{}BGE-M3 & 0.5620 & 0.7214 & +\allowbreak{}0.1594 & +\allowbreak{}28.3680\% \\
Zalo B & BM25+\allowbreak{}BGE-M3+\allowbreak{}E5 & 0.5620 & 0.7344 & +\allowbreak{}0.1724 & +\allowbreak{}30.6830\% \\
BCA combined descriptive & BM25+\allowbreak{}BGE-M3 & 0.1687 & 0.2283 & +\allowbreak{}0.0597 & +\allowbreak{}35.3820\% \\
BCA combined descriptive & BM25+\allowbreak{}BGE-M3+\allowbreak{}E5 & 0.1687 & 0.2297 & +\allowbreak{}0.0610 & +\allowbreak{}36.1760\% \\
\end{longtable}
\endgroup

The hybrid gains over BM25 are large in Zalo and smaller but still positive in ALQAC. BCA also shows a positive descriptive difference from BM25, but its lower absolute scores must be interpreted with coverage limits and component dependence. The answer to RQ1 is therefore positive within the evaluated conditions: lexical–semantic hybrids improve the fixed BM25 reference. The result should not be generalized into a claim that one hybrid configuration will transfer unchanged to every legal corpus.

\section{Incremental Contribution of E5}

This section tests whether E5 adds value beyond BM25+BGE-M3. This is a stricter question than asking whether E5 is a competent standalone retriever. The relevant comparison is the two-way hybrid against the three-way hybrid.

\begingroup
\footnotesize
\setlength{\tabcolsep}{3pt}
\renewcommand{\arraystretch}{1.18}
\begin{longtable}{>{\raggedright\arraybackslash}p{2.272cm}>{\raggedright\arraybackslash}p{2.272cm}>{\raggedright\arraybackslash}p{2.429cm}>{\raggedright\arraybackslash}p{2.515cm}>{\raggedright\arraybackslash}p{2.515cm}>{\raggedright\arraybackslash}p{2.677cm}}
\caption[Incremental E5 contribution on the canonical ALQAC and Zalo populations.]{Incremental E5 contribution on the canonical ALQAC and Zalo populations.}\label{tab:5.4}\\
\toprule
\textbf{Condition} & \textbf{Two-way NDCG} & \textbf{Three-way NDCG} & \textbf{Absolute change} & \textbf{Relative change} & \textbf{Stored Wilcoxon p} \\
\midrule
\endfirsthead
\toprule
\textbf{Condition} & \textbf{Two-way NDCG} & \textbf{Three-way NDCG} & \textbf{Absolute change} & \textbf{Relative change} & \textbf{Stored Wilcoxon p} \\
\midrule
\endhead
\midrule
\multicolumn{6}{r}{\footnotesize Continued on next page}\\
\endfoot
\bottomrule
\endlastfoot
ALQAC A & 0.8805 & 0.8784 & -0.0021 & -0.2380\% & 0.4498 \\
ALQAC B & 0.7067 & 0.7055 & -0.0012 & -0.1630\% & 0.7449 \\
Zalo A & 0.6805 & 0.6818 & +\allowbreak{}0.0013 & +\allowbreak{}0.1900\% & 0.1548 \\
Zalo B & 0.7214 & 0.7344 & +\allowbreak{}0.0130 & +\allowbreak{}1.8030\% & \(5.1598\times10^{-7}\) \\
\end{longtable}
\endgroup

The incremental E5 result is condition-dependent. ALQAC shows small negative NDCG@10 differences in both corpus representations. Zalo A shows a very small positive difference, while Zalo B shows the clearest positive result: an absolute NDCG@10 gain of approximately 0.0130, with higher MRR@10 and Hit@10 as well. The p-values in Table~\ref{tab:5.4} are reported as stored query-level paired tests; they do not remove the need to interpret multiple conditions and possible query dependence carefully.

The A/B corpus comparison also resists a simple parser narrative. ALQAC has substantially lower hybrid scores in the parsed condition than in the clean condition. Zalo, in contrast, has higher scores in the parsed-with-fallback condition than in the clean condition. This does not imply that parsing is inherently harmful for ALQAC and inherently beneficial for Zalo. The two collections differ in construction, target mapping, fallback behavior, and article universe. A matched no-fallback ablation was not available, so the independent effect of fallback remains unmeasured.

The answer to RQ2 is therefore not a general yes. E5 provides useful incremental value in the parsed Zalo condition, negligible or unsupported value in clean Zalo, and no observed improvement in ALQAC. This is one of the thesis's main practical findings: adding a dense retriever expands the search space, but it does not guarantee better ranking on unseen data.

\section{BCA Coverage and Component-Aware Evaluation}

BCA serves two purposes. It extends the evaluation beyond ALQAC and Zalo, and it exposes how corpus coverage and target dependence affect interpretation. Table~\ref{tab:5.5} reports BCA results under the giant-component, small-component, and combined descriptive views.

\begin{table}[H]
\centering
\scriptsize
\setstretch{1.0}
\setlength{\tabcolsep}{3pt}
\renewcommand{\arraystretch}{0.93}
\caption[Compact BCA protocol and coverage views.]{Compact BCA protocol and coverage views. All metrics are at cutoff 10; OOF denotes out-of-fold.}\label{tab:5.5}
\vspace{6pt}
\resizebox{\textwidth}{!}{%
\begin{tabular}{@{}lllrrrrr@{}}
\toprule
\textbf{Protocol} & \textbf{Cohort} & \textbf{Method} & \textbf{N} & \textbf{NDCG} & \textbf{MRR} & \textbf{MAP} & \textbf{Hit} \\
\midrule
Giant & All & BM25 & 484 & 0.1670 & 0.2469 & 0.1098 & 0.4545 \\
Giant & All & BGE-M3 & 484 & 0.2153 & 0.3143 & 0.1453 & 0.5455 \\
Giant & All & BM25+BGE & 484 & 0.2275 & 0.3414 & 0.1573 & 0.5517 \\
Giant & All & BM25+BGE+E5 & 484 & 0.2295 & 0.3413 & 0.1560 & 0.5702 \\
Giant & Answerable & BM25 & 394 & 0.2051 & 0.3033 & 0.1348 & 0.5584 \\
Giant & Answerable & BGE-M3 & 394 & 0.2645 & 0.3861 & 0.1785 & 0.6701 \\
Giant & Answerable & BM25+BGE & 394 & 0.2795 & 0.4194 & 0.1932 & 0.6777 \\
Giant & Answerable & BM25+BGE+E5 & 394 & 0.2819 & 0.4193 & 0.1916 & 0.7005 \\
\midrule
Small OOF & All & BM25 & 62 & 0.1817 & 0.2157 & 0.1382 & 0.3387 \\
Small OOF & All & BGE-M3 & 62 & 0.2015 & 0.2392 & 0.1597 & 0.3387 \\
Small OOF & All & BM25+BGE & 62 & 0.2346 & 0.2921 & 0.1865 & 0.3871 \\
Small OOF & All & BM25+BGE+E5 & 62 & 0.2309 & 0.2838 & 0.1815 & 0.3871 \\
Small OOF & Answerable & BM25 & 46 & 0.2449 & 0.2907 & 0.1863 & 0.4565 \\
Small OOF & Answerable & BGE-M3 & 46 & 0.2715 & 0.3225 & 0.2153 & 0.4565 \\
Small OOF & Answerable & BM25+BGE & 46 & 0.3162 & 0.3937 & 0.2514 & 0.5217 \\
Small OOF & Answerable & BM25+BGE+E5 & 46 & 0.3112 & 0.3825 & 0.2446 & 0.5217 \\
\midrule
Combined OOF & All & BM25 & 546 & 0.1687 & 0.2434 & 0.1130 & 0.4414 \\
Combined OOF & All & BGE-M3 & 546 & 0.2137 & 0.3058 & 0.1469 & 0.5220 \\
Combined OOF & All & BM25+BGE & 546 & 0.2283 & 0.3358 & 0.1606 & 0.5330 \\
Combined OOF & All & BM25+BGE+E5 & 546 & 0.2297 & 0.3348 & 0.1589 & 0.5495 \\
Combined OOF & Answerable & BM25 & 440 & 0.2093 & 0.3020 & 0.1402 & 0.5477 \\
Combined OOF & Answerable & BGE-M3 & 440 & 0.2652 & 0.3795 & 0.1823 & 0.6477 \\
Combined OOF & Answerable & BM25+BGE & 440 & 0.2833 & 0.4167 & 0.1993 & 0.6614 \\
Combined OOF & Answerable & BM25+BGE+E5 & 440 & 0.2850 & 0.4154 & 0.1972 & 0.6818 \\
\bottomrule
\end{tabular}}
\end{table}

The BCA reference graph contains one giant component with 484 of 546 queries and 49 smaller components containing the remaining 62 queries. This structure affects statistical interpretation. The combined descriptive view shows the three-way hybrid slightly above the two-way hybrid in NDCG@10, but the small-component out-of-fold view shows a negative NDCG@10 difference for the three-way system. The corresponding component-aware inference does not establish an E5 advantage.

Coverage is also a major limitation. Of the 546 BCA queries, 440 have at least one represented target and 106 have none. A further 229 answerable queries are partial-target cases in which some, but not all, targets are present. An unanswerable query cannot receive a positive Hit@10 from any ranking method over the fixed corpus. Therefore, BCA's lower end-to-end scores partly reflect corpus and mapping limitations rather than only scoring failure.

The BCA result is a dataset and evaluation contribution rather than a simple model leaderboard. It shows that an end-to-end legal retrieval score can mix three phenomena: whether the evidence exists, whether it is mapped to the correct legal identity, and whether the retriever ranks it highly. A method improvement can address only the third problem unless corpus construction changes as well.

\section{Failure Analysis}

The failure analysis asks what a model intervention could plausibly repair. Missing targets require corpus or mapping work. Available targets outside the top ten require ranking or fusion improvements. Relevant articles demoted within the top ten affect NDCG and MRR even when Hit@10 is unchanged. These distinctions guide the adaptation experiment but do not prove that fine-tuning will solve the observed failures.

BCA provides the clearest evidence of non-rankable failures: 106 questions have no mapped target in the corpus. For these questions, changing BM25 weights, dense encoders, or fusion rules cannot recover the target. Partial-target queries create a second limitation: a system may retrieve an available relevant article while still being penalized for missing other targets that are absent from the corpus.

For ALQAC and Zalo, top-ten boundary cases show that E5 can both recover and lose hits. The same added signal that improves parsed Zalo on average can demote specific relevant articles below the cutoff. Appendix E preserves representative cases. The main conclusion is more general: aggregate improvement does not imply uniform query-level improvement, and aggregate regression does not imply that every query becomes worse.

The analysis remains an information-retrieval analysis. It identifies target absence, top-ten transitions, within-law article demotion, and persistent misses. It does not assign legal-reasoning causes such as misunderstanding exceptions or temporal validity. A legally adjudicated error taxonomy would require a sampling plan, full text review, annotation rules, and reviewer agreement. That study was not performed here.

\section{Legal-Domain Adaptation}

This section examines whether legal-domain BGE-M3 adaptation improves retrieval beyond development selection. The narrative has six parts: Study A development selection, Study A corrected hybrid evaluation on unseen data, Study B dense-only adaptation, transfer of dense improvements into hybrid retrieval, query-level analysis, and overall interpretation.

\subsection{Development Selection}

Study A: Hybrid-Oriented Adaptation selects a BGE-M3 training recipe on development queries. Table~\ref{tab:5.7} reports the learning-rate and epoch pilot under passage length 512. Table~\ref{tab:5.8} reports the subsequent batch ablation at the selected learning rate and epoch.

\begingroup
\small
\setlength{\tabcolsep}{3pt}
\renewcommand{\arraystretch}{1.18}
\begin{longtable}{>{\raggedright\arraybackslash}p{3.672cm}>{\raggedright\arraybackslash}p{3.564cm}>{\raggedright\arraybackslash}p{3.810cm}>{\raggedright\arraybackslash}p{4.073cm}}
\caption[Passage-512 learning-rate and epoch pilot: development NDCG@10.]{Passage-512 learning-rate and epoch pilot: development NDCG@10.}\label{tab:5.7}\\
\toprule
\textbf{Learning rate} & \textbf{Epoch} & \textbf{Two-way hybrid} & \textbf{Three-way hybrid} \\
\midrule
\endfirsthead
\toprule
\textbf{Learning rate} & \textbf{Epoch} & \textbf{Two-way hybrid} & \textbf{Three-way hybrid} \\
\midrule
\endhead
\midrule
\multicolumn{4}{r}{\footnotesize Continued on next page}\\
\endfoot
\bottomrule
\endlastfoot
\(5.0000\times10^{-6}\) & 1 & 0.7470 & 0.7685 \\
\(5.0000\times10^{-6}\) & 2 & 0.7603 & 0.7785 \\
\(5.0000\times10^{-6}\) & 3 & 0.7647 & 0.7815 \\
\(1.0000\times10^{-5}\) & 1 & 0.7509 & 0.7816 \\
\(1.0000\times10^{-5}\) & 2 & 0.7658 & 0.7873 \\
\(1.0000\times10^{-5}\) & 3 & 0.7628 & 0.7897 \\
\(2.0000\times10^{-5}\) & 1 & 0.7440 & 0.7757 \\
\(2.0000\times10^{-5}\) & 2 & 0.7691 & 0.7923 \\
\(2.0000\times10^{-5}\) & 3 & 0.7689 & 0.7958 \\
\end{longtable}
\endgroup

The pilot shows development improvements across several candidate settings. The selected recipe uses learning rate 0.00002 and two epochs. The choice is not evidence that the selected point is a global optimum; it is the best supported point under the restricted search and selection rule.

\begin{table}[H]
\centering
\footnotesize
\setstretch{1.0}
\setlength{\tabcolsep}{6pt}
\renewcommand{\arraystretch}{1.05}
\caption[Batch ablation at the selected learning rate and epoch.]{Batch ablation at the selected learning rate and epoch; scores are development NDCG@10.}\label{tab:5.8}
\vspace{6pt}
\begin{tabular}{@{}lrrrr@{}}
\toprule
\textbf{Batch/accum.} & \textbf{Two-way} & \textbf{$\Delta$} & \textbf{Three-way} & \textbf{$\Delta$} \\
\midrule
Pretrained baseline & 0.7172 & -- & 0.7245 & -- \\
1 / 32 & 0.7725 & +0.0553 & 0.7869 & +0.0624 \\
2 / 16 & 0.7693 & +0.0521 & 0.7924 & +0.0679 \\
4 / 8 & 0.7781 & +0.0609 & 0.7934 & +0.0689 \\
\bottomrule
\end{tabular}
\end{table}

The selected development configuration improves two-way NDCG@10 by 0.060945 and three-way NDCG@10 by 0.068882. This justifies testing the intervention on unseen data. It does not establish generalization, because the development queries are part of model selection.

\subsection{Corrected Hybrid Evaluation on Unseen Data}

The corrected comparison on unseen data evaluates Zalo parsed fold 0 with 640 queries. Fusion weights are fitted on the permitted training-side folds using pretrained BGE-M3 and then held fixed when BGE-M3 is replaced by its adapted version. This design measures encoder replacement within a specified hybrid rather than unrestricted retuning of the whole system.

The corrected two-way weights are BM25/BGE/E5 = 0.2/0.8/0.0. The corrected three-way weights are 0.2/0.5/0.3. Because a flawed preliminary analysis had already exposed the same fold, the corrected comparison is interpreted as a post-hoc evaluation on data unseen during training and selection, not as a pristine first-use test.

\begin{table}[H]
\centering
\footnotesize
\setstretch{1.0}
\setlength{\tabcolsep}{5pt}
\renewcommand{\arraystretch}{1.05}
\caption[Corrected Zalo fold-0 two-way hybrid evaluation.]{Corrected Zalo fold-0 two-way hybrid evaluation; all metrics are at cutoff 10.}\label{tab:5.9}
\vspace{6pt}
\begin{tabular}{@{}llrrrrr@{}}
\toprule
\textbf{BGE state} & \textbf{Seed} & \textbf{NDCG} & \textbf{$\Delta$} & \textbf{MRR} & \textbf{MAP} & \textbf{Hit} \\
\midrule
Pretrained & Fixed & 0.7516 & -- & 0.7015 & 0.6985 & 0.9141 \\
Fine-tuned & 42 & 0.7005 & \ensuremath{-}0.0511 & 0.6470 & 0.6451 & 0.8750 \\
Fine-tuned & 43 & 0.7005 & \ensuremath{-}0.0511 & 0.6490 & 0.6467 & 0.8688 \\
Fine-tuned & 44 & 0.6977 & \ensuremath{-}0.0539 & 0.6398 & 0.6374 & 0.8859 \\
Fine-tuned & Mean & 0.6995 & \ensuremath{-}0.0520 & 0.6453 & 0.6430 & 0.8766 \\
\bottomrule
\end{tabular}
\end{table}

All three two-way seeds score below the pretrained two-way baseline. The mean NDCG@10 change is -0.052045, and MRR@10, MAP@10, and Hit@10 are also lower. The result therefore cannot be rescued by choosing a different seed from the three reported runs.

\begin{table}[H]
\centering
\footnotesize
\setstretch{1.0}
\setlength{\tabcolsep}{5pt}
\renewcommand{\arraystretch}{1.05}
\caption[Corrected Zalo fold-0 three-way hybrid evaluation.]{Corrected Zalo fold-0 three-way hybrid evaluation; all metrics are at cutoff 10.}\label{tab:5.10}
\vspace{6pt}
\begin{tabular}{@{}llrrrrr@{}}
\toprule
\textbf{BGE state} & \textbf{Seed} & \textbf{NDCG} & \textbf{$\Delta$} & \textbf{MRR} & \textbf{MAP} & \textbf{Hit} \\
\midrule
Pretrained & Fixed & 0.7595 & -- & 0.7077 & 0.7048 & 0.9266 \\
Fine-tuned & 42 & 0.7320 & \ensuremath{-}0.0275 & 0.6771 & 0.6743 & 0.9141 \\
Fine-tuned & 43 & 0.7340 & \ensuremath{-}0.0255 & 0.6771 & 0.6747 & 0.9203 \\
Fine-tuned & 44 & 0.7378 & \ensuremath{-}0.0217 & 0.6794 & 0.6771 & 0.9281 \\
Fine-tuned & Mean & 0.7346 & \ensuremath{-}0.0249 & 0.6779 & 0.6754 & 0.9208 \\
\bottomrule
\end{tabular}
\end{table}

The three-way hybrid also has lower mean NDCG@10 after fine-tuning. The regression is smaller than in the two-way system, but it remains a regression against the matched pretrained three-way comparator. Seed 44 slightly improves Hit@10 while reducing NDCG@10 and MRR@10, showing why the primary rank-sensitive metric must remain central.

\begin{table}[H]
\centering
\footnotesize
\setstretch{1.0}
\setlength{\tabcolsep}{4pt}
\renewcommand{\arraystretch}{1.05}
\caption[Corrected paired statistics for NDCG@10.]{Corrected paired statistics for NDCG@10.}\label{tab:5.11}
\vspace{6pt}
\resizebox{\textwidth}{!}{%
\begin{tabular}{@{}lrrrrrr@{}}
\toprule
\textbf{Comparison} & \textbf{Queries} & \textbf{Clusters} & \textbf{Mean $\Delta$} & \textbf{95\% cluster CI} & \textbf{Raw p} & \textbf{Holm p} \\
\midrule
Fine-tuned minus pretrained, two-way & 640 & 106 & \ensuremath{-}0.0520 & [\ensuremath{-}0.0919, +\allowbreak{}0.1076] & 0.7480 & 1.0000 \\
Fine-tuned minus pretrained, three-way & 640 & 106 & \ensuremath{-}0.0249 & [\ensuremath{-}0.0542, +\allowbreak{}0.0938] & 0.7480 & 1.0000 \\
\bottomrule
\end{tabular}}
\end{table}

The intervals in Table~\ref{tab:5.11} include zero, so the observed regressions are not described as statistically established population-level harm. However, non-significance is not equivalence, and the practical improvement threshold of +0.01 NDCG@10 is not met. The correct conclusion is that the selected adaptation recipe failed to demonstrate hybrid improvement on unseen data under this protocol.

\subsection{Study B: Dense-Oriented Adaptation}

Study A does not by itself determine whether adaptation failed at the dense-retriever level or only after integration into the hybrid. Study B directly evaluates dense-only BGE-M3 adaptation before testing hybrid transfer.

\begingroup
\footnotesize
\setlength{\tabcolsep}{3pt}
\renewcommand{\arraystretch}{1.18}
\begin{longtable}{>{\raggedright\arraybackslash}p{3.587cm}>{\raggedright\arraybackslash}p{2.807cm}>{\raggedright\arraybackslash}p{2.892cm}>{\raggedright\arraybackslash}p{2.807cm}>{\raggedright\arraybackslash}p{2.807cm}}
\caption[Study B post-hoc dense benchmark on 640 official queries.]{Study B post-hoc dense benchmark on 640 official queries.}\label{tab:5.12}\\
\toprule
\textbf{System} & \textbf{NDCG@10} & \textbf{Delta NDCG@10} & \textbf{MRR@10} & \textbf{Recall@100} \\
\midrule
\endfirsthead
\toprule
\textbf{System} & \textbf{NDCG@10} & \textbf{Delta NDCG@10} & \textbf{MRR@10} & \textbf{Recall@100} \\
\midrule
\endhead
\midrule
\multicolumn{5}{r}{\footnotesize Continued on next page}\\
\endfoot
\bottomrule
\endlastfoot
Pretrained & 0.6740 & 0.0000 & 0.6213 & 0.9617 \\
Fine-tuned seed 42 & 0.6932 & 0.0192 & 0.6448 & 0.9531 \\
Fine-tuned seed 43 & 0.6925 & 0.0185 & 0.6431 & 0.9516 \\
Fine-tuned seed 44 & 0.6928 & 0.0189 & 0.6437 & 0.9531 \\
Fine-tuned seed mean & 0.6928 & 0.0189 & 0.6438 & 0.9526 \\
\end{longtable}
\endgroup

Study B shows a mean NDCG@10 gain of 0.018866, with all three seeds improving this endpoint. Mean Recall@100 decreases by 0.009115. The dense result therefore supports the claim that legal-domain fine-tuning can improve the dense encoder on this fixed-corpus query benchmark, while also showing a small recall tradeoff at depth 100.

\subsection{Transfer of Dense Improvements into Hybrid Retrieval}

The crucial question is whether the dense improvement remains useful inside the full hybrid. Table~\ref{tab:5.13} reports the Study B primary locked-weight hybrid outcomes. Fusion weights are selected with pretrained BGE-M3 and then held fixed for the corresponding fine-tuned comparison.

\begin{table}[H]
\centering
\footnotesize
\setstretch{1.0}
\setlength{\tabcolsep}{4pt}
\renewcommand{\arraystretch}{1.05}
\caption[Study B primary fusion outcomes under pretrained-locked weights.]{Study B primary fusion outcomes under pretrained-locked weights.}\label{tab:5.13}
\vspace{6pt}
\resizebox{\textwidth}{!}{%
\begin{tabular}{@{}lrrrrrrl@{}}
\toprule
\textbf{Hybrid} & \textbf{Pretrained} & \textbf{FT mean} & \textbf{$\Delta$ NDCG} & \textbf{Rel. $\Delta$} & \textbf{$\Delta$ R@100} & \textbf{Positive seeds} & \textbf{Gate} \\
\midrule
Two-way & 0.7526 & 0.7483 & -0.0043 & -0.5707\% & -0.0057 & 0/\allowbreak{}3 & Fail \\
Three-way & 0.7467 & 0.7476 & 0.0009 & 0.1203\% & -0.0031 & 2/\allowbreak{}3 & Fail \\
\bottomrule
\end{tabular}}
\end{table}

The dense-only gain does not transfer into a practically useful hybrid gain. The two-way hybrid decreases by 0.004295 NDCG@10. The three-way hybrid increases by only 0.000898 NDCG@10, far below the +0.01 improvement threshold. Thus, Study B sharpens the main RQ3 finding: adaptation can improve standalone dense ranking while failing to improve the complete hybrid retrieval system by a meaningful amount.

The mechanism is not isolated by these experiments. Possible explanations include changed score calibration, overlap between lexical and dense evidence, and fusion weights selected for pretrained scores. Because these mechanisms are not independently varied, they remain hypotheses for future work rather than established causes.

\subsection{Query-Level Analysis}

Aggregate means can hide mixtures of repaired and introduced failures. Table~\ref{tab:5.14} summarizes the Study A corrected seed-42 hybrid changes. The first three rows partition all 640 queries by NDCG@10 change; the final rows describe top-ten transitions and are not additional disjoint totals.

\begingroup
\small
\setlength{\tabcolsep}{3pt}
\renewcommand{\arraystretch}{1.18}
\begin{longtable}{>{\raggedright\arraybackslash}p{6.897cm}>{\raggedright\arraybackslash}p{4.080cm}>{\raggedright\arraybackslash}p{4.362cm}}
\caption[Query-level changes after fine-tuning, corrected seed 42.]{Query-level changes after fine-tuning, corrected seed 42.}\label{tab:5.14}\\
\toprule
\textbf{Outcome} & \textbf{Two-way hybrid} & \textbf{Three-way hybrid} \\
\midrule
\endfirsthead
\toprule
\textbf{Outcome} & \textbf{Two-way hybrid} & \textbf{Three-way hybrid} \\
\midrule
\endhead
\midrule
\multicolumn{3}{r}{\footnotesize Continued on next page}\\
\endfoot
\bottomrule
\endlastfoot
Higher NDCG@10 & 102 & 102 \\
Lower NDCG@10 & 210 & 163 \\
Unchanged NDCG@10 & 328 & 375 \\
Previously missing at ten, now retrieved & 23 & 16 \\
Previously retrieved at ten, now missing & 48 & 24 \\
Missing at ten in both systems & 32 & 31 \\
\end{longtable}
\endgroup

The two-way system repairs some top-ten misses, but introduces more than twice as many new top-ten misses as it repairs. The three-way system introduces fewer new misses, consistent with its smaller aggregate regression, but it still introduces more misses than it repairs. These counts show why the negative mean should not be interpreted as uniform deterioration, and why isolated repaired examples cannot overturn the aggregate result.

\begingroup
\small
\setlength{\tabcolsep}{3pt}
\renewcommand{\arraystretch}{1.18}
\begin{longtable}{>{\raggedright\arraybackslash}p{3.780cm}>{\raggedright\arraybackslash}p{3.780cm}>{\raggedright\arraybackslash}p{3.780cm}>{\raggedright\arraybackslash}p{3.780cm}}
\caption[Study B NDCG changes by query, using mean fine-tuned seed metric minus pretrained.]{Study B NDCG changes by query, using mean fine-tuned seed metric minus pretrained.}\label{tab:5.15}\\
\toprule
\textbf{Comparison} & \textbf{Improved} & \textbf{Worse} & \textbf{Tied} \\
\midrule
\endfirsthead
\toprule
\textbf{Comparison} & \textbf{Improved} & \textbf{Worse} & \textbf{Tied} \\
\midrule
\endhead
\midrule
\multicolumn{4}{r}{\footnotesize Continued on next page}\\
\endfoot
\bottomrule
\endlastfoot
Dense & 127 & 96 & 417 \\
Two-way & 81 & 93 & 466 \\
Three-way & 85 & 86 & 469 \\
\end{longtable}
\endgroup

Study B's dense-only comparison has more improved than worsened queries, while the Study B hybrids do not show the same favorable balance. This query-level view is consistent with the aggregate finding: dense improvement exists, but much of it is absorbed or neutralized when the adapted encoder is placed inside the full fusion system.

\subsection{Overall Interpretation}

The RQ3 evidence has a clear sequence. First, Study A's development selection shows large gains and justifies testing the adaptation. Second, Study A's corrected hybrid evaluation on unseen data shows regressions for both two-way and three-way hybrids. Third, Study B shows that legal fine-tuning can improve dense BGE-M3 ranking. Fourth, that dense gain largely fails to transfer into locked-weight hybrid retrieval.

This sequence supports a stronger and more nuanced conclusion than either ``fine-tuning works'' or ``fine-tuning fails.'' It works at the dense-retriever level in Study B, but it does not necessarily improve the complete lexical–semantic system. The complete system must therefore be evaluated directly. Development scores and dense-only scores are insufficient substitutes for end-to-end hybrid evaluation.

Several boundaries remain. Cross-collection adaptation to ALQAC and BCA was not established. Secondary retuned-weight hybrid results were not available, so the thesis does not claim that the adapted encoder could not be made useful under a different fusion policy. A fresh untouched cohort would be required for a cleaner confirmatory test. These limitations narrow the conclusion without weakening the observed negative transfer result.

\section{Experimental Synthesis}

The experiments answer the research questions as follows. For RQ1, hybrid retrieval improves the fixed BM25 reference in the evaluated ALQAC, Zalo, and descriptive BCA conditions. For RQ2, E5 has condition-dependent incremental value: parsed Zalo provides the clearest positive case, while ALQAC and BCA do not support a general three-way advantage. For RQ3, legal-domain adaptation improves development scores and can improve dense-only ranking, but it does not necessarily improve the complete hybrid retrieval system.

The BCA analysis adds a methodological lesson. Evaluation cannot be reduced to a ranking table when target coverage and query dependence are substantial. Missing targets, partial target availability, and giant connected components affect what a retrieval score means. Reporting both end-to-end and answerable-only views makes this limitation visible rather than hiding it inside an average.

The principal scientific takeaway is that legal retrieval systems must be evaluated at the level at which they will be used. A model can be better in isolation while adding little to a hybrid; an added retriever can help one corpus representation and not another; a retrieval failure can be caused by absent evidence rather than poor ranking. These findings justify the thesis's emphasis on controlled comparison, negative results, and careful limitation of causal claims.
